\documentclass[letterpaper]{article}
\usepackage{../../common/fspdoc}
\setcounter{secnumdepth}{5}
\setcounter{tocdepth}{5}




% Recipe commands to run initially only in CI

% begin_fsp_run 
% fsp_validation_comment - Remove WW file update crontab entry (CI only command)
% [master]$ rm -f /etc/cron.d/wwupdatefiles
% [master]$ TIME_T1=$(date +%s)
% [master]$ echo " "; echo "BEGIN TIMING"
% end_fsp_run

% Include git variables
%%% This file has been generated by the vc bundle for TeX.
%%% Do not edit this file!
%%%
%%% Define Git specific macros.
\gdef\GITHash{b1e5ea39b0eb099e892994a2f56798d77c1b41ad}%
\gdef\GITAbrHash{b1e5ea3}%
\gdef\GITParentHashes{ab0141b6f1cd0193472944a1642275fb792a8f77}%
\gdef\GITAbrParentHashes{ab0141b}%
\gdef\GITAuthorName{Karl W. Schulz}%
\gdef\GITAuthorEmail{karl.w.schulz@intel.com}%
\gdef\GITAuthorDate{2015-04-17 16:52:14 -0500}%
\gdef\GITCommitterName{Karl W. Schulz}%
\gdef\GITCommitterEmail{karl.w.schulz@intel.com}%
\gdef\GITCommitterDate{2015-04-17 16:52:14 -0500}%
%%% Define generic version control macros.
\gdef\VCRevision{\GITAbrHash}%
\gdef\VCAuthor{\GITAuthorName}%
\gdef\VCDateRAW{2015-04-17}%
\gdef\VCDateISO{2015-04-17}%
\gdef\VCDateTEX{2015/04/17}%
\gdef\VCTime{16:52:14 -0500}%
\gdef\VCModifiedText{\textcolor{red}{with local modifications!}}%
%%% Assume clean working copy.
\gdef\VCModified{0}%
\gdef\VCRevisionMod{\VCRevision}%


% Define Base OS and other local macros
\newcommand{\baseOS}{CentOS7.1}
\newcommand{\LosF}{\emph{LosF*}}
\newcommand{\FSPVersion}{15.27}

\lhead{{2015 \FSP{} Install Guide - {\small \em \baseOS{} Version}
    (v\FSPVersion{}) } }
%\fancyfoot[RE,RO]{\small Last Update: \VCDateISO~(Rev: \VCRevision) - v\FSPVersion{}}

\newcommand{\install}{yum -y install}
\newcommand{\chrootinstall}{yum -y --installroot=\$CHROOT install}
\newcommand{\groupinstall}{yum -y groupinstall}
\newcommand{\groupchrootinstall}{yum -y --installroot=\$CHROOT groupinstall}

\begin{document}
\graphicspath{{../../common/figures/}}
\thispagestyle{empty}

% Title Page

{\hspace*{4in} \includegraphics[width=1.8in]{intel_1spot_100.pdf}}

\vspace*{2cm}
\noindent {\LARGE \color{RoyalBlue} \fontfamily{phv}\selectfont 2015 Forest Peak (v\FSPVersion{})} \vspace*{0.1cm} \\
\noindent {\LARGE \color{RoyalBlue} \fontfamily{phv}\selectfont Cluster Building Recipes} \\ 
\noindent\rule{6in}{0.4pt} \\

\noindent {\Large \baseOS{} Base OS} \\ \vspace{0.2cm}

\noindent {\large {\em Base Linux* Edition }}

\vspace*{3in}

\noindent{\normalsize Intel Cluster Makers} \vspace*{0.1cm} \\
{\normalsize Copyright~{\small\copyright}~2014-2015 Intel Corporation} \vspace*{0.1cm} \\ 
{\normalsize Document Last Update: \VCDateISO} \vspace*{0.1cm} \\ 
{\normalsize Document Revision: \VCRevision} \\ \vspace*{0.1cm}

% Disclaimer  ----------------------------------------------------
No license (express or implied, by estoppel or otherwise) to any intellectual property rights is granted by this document. 
\newline \\
\noindent This document contains information on products, services and/or processes in development.  All information provided here is subject to change without notice. Contact your Intel representative to obtain the latest forecast, schedule, specifications and roadmaps.
\newline \\
\noindent The products and services described may contain defects or errors known as errata which may cause deviations from published specifications. Current characterized errata are available on request.
\newline \\
\noindent You may not use or facilitate the use of this document in connection with any infringement or other legal analysis concerning Intel products described herein. You agree to grant Intel a non-exclusive, royalty-free license to any patent claim thereafter drafted which includes subject matter disclosed herein.
\newline \\
\noindent Forecasts:  Any forecasts of requirements for goods and services are provided for discussion purposes only. Intel will have no liability to make any purchase pursuant to forecasts.  Any cost or expense you incur to respond to requests for information or in reliance on any forecast will be at your own risk and expense.
\newline \\
\noindent Copies of documents which have an order number and are referenced in this document may be obtained by calling 1-800-548-4725 or by visiting www.intel.com/design/literature.htm.
\newline \\
\noindent Intel technologies’ features and benefits depend on system configuration and may require enabled hardware, software or service activation. Learn more at Intel.com, or from the OEM or retailer.
\newline \\
\noindent Software and workloads used in performance tests may have been optimized for performance only on Intel microprocessors.
\newline \\
\noindent Performance tests, such as SYSmark and MobileMark, are measured using specific computer systems, components, software, operations and functions. Any change to any of those factors may cause the results to vary. You should consult other information and performance tests to assist you in fully evaluating your contemplated purchases, including the performance of that product when combined with other products.
\newline \\
\noindent For more complete information visit http://www.intel.com/performance.  
\newline \\
\noindent Intel, the Intel logo, \Intel{}~MPI, \Intel{} Cluster Checker, \FSP{} are trademarks of Intel Corporation in the U.S. and/or other countries.
\newline \\
\noindent *Other names and brands may be claimed as the property of others.  
\newline \\
\noindent {\small\copyright} 2015 Intel Corporation

 

\newpage
\tableofcontents
\newpage

% Introduction  --------------------------------------------------

\section{Introduction} \label{sec:introduction}
This guide presents a simple cluster installation procedure using components
from the Forest Peak (\FSP{}) software stack. \FSP{} represents an aggregation of a number of
common ingredients required to deploy and manage an HPC Linux* cluster including
provisioning tools, resource management, I/O clients, development tools, and a variety of
scientific libraries. These packages have been pre-built with HPC integration
in mind and represent a mix of open-source components combined with \Intel{}
development and analysis tools (e.g. \Intel{} Parallel Studio XE Cluster Edition).
%This install guide assumes availablility of a {\em master} host
%that will have access to common OS repositories and the \FSP{} repository in order
%to facilitate software installs and dependency resolution.  
The documentation herein is intended to be reasonably generic, but uses the
underlying motivation of a small, 4-node cluster install to define a step-by-step
process. Several optional customizations are included and the intent is that
these collective instructions can be modified as needed for local site
customizations. 
 \\

\noindent {\bf Base Linux Edition}: this edition of the guide highlights
installation without the use of a companion configuration management system and
directly uses distro-provided package management tools for component
selection. The steps that follow also highlight specific changes to system
configuration files that are required as part of the cluster install
process. 
%Other editions of this guide provide similar install steps when using
%specific configuration management systems that can simplify the installation
%and configuration process. 

\subsection{Target Audience}

This guide is targeted at experienced \Linux{} system administrators for HPC
environments. Knowledge of software package management, system networking, and
PXE booting is assumed. Command-line input examples are highlighted throughout
this guide via the following syntax:

\begin{lstlisting}[language=bash,literate={-}{-}1,keywords={},upquote=true]
[sms](*\#*) echo "OpenHPC hello world"
\end{lstlisting}

Unless specified otherwise, the examples presented are executed with
elevated (root) privileges. The examples also presume use of the BASH login
shell, though the equivalent commands in other shells can be substituted.
In addition to specific command-line instructions called out in this guide, an
alternate convention is used to highlight potentially useful tips or optional
configuration options. These tips are highlighted via the following format:

\begin{center}
\begin{tcolorbox}[]
\small
I propose we leave math to the machines and go play outside. --Calvin
\end{tcolorbox}
\end{center}


\noindent {\bf Requirements}: this recipe assumes the availability of a single head node
{\em master}, and four {\em compute} nodes. The {\em master} node is
provisioned with \baseOS{} and is subsequently configured to
provision the remaining {\em compute} nodes with Warewulf in a
stateless configuration. We assume an external \$HOME file system is
available for use on both the master and compute nodes via NFS, and that compute node BMCs are
available via IPMI from the chosen master host. \\



% -*- mode: latex; fill-column: 120; -*- 

\subsection{Inputs} \label{sec:inputs}
As this recipe details installing a cluster starting from bare-metal, there is a requirement to define IP addresses and
gather hardware MAC addresses in order to support a controlled provisioning process. These values are necessarily unique
to the hardware being used, and this document uses variable substitution (\texttt{\$\{variable\}}) in the command-line
examples that follow to highlight where local site inputs are required. A summary of the required and optional variables
used throughout this recipe are presented below. Note that while the example definitions above correspond to a small
4-node compute subsystem, the compute parameters are defined in array format to accommodate logical extension to larger
node counts. \\

\vspace*{0.2cm}
\begin{tabular}{@{}>{\textbullet}l p{7cm} l}
& \texttt{\$\{ohpc\_repo\}} & {\small \# OpenHPC repo location} \\
& \texttt{\$\{master\_name\}} & {\small \# Hostname for master server} \\
& \texttt{\$\{master\_ip\}} & {\small \# Internal IP address on master server}  \\
& \texttt{\$\{sms\_eth\_internal\}} & {\small \# Internal Ethernet interface on SMS} \\
& \texttt{\$\{eth\_provision\}} & {\small \# Provisioning interface for computes} \\
& \texttt{\$\{internal\_netmask\}} & {\small \# Subnet netmask for internal network} \\
& \texttt{\$\{ntp\_server\}} & {\small \# Local ntp server for time synchronization} \\
& \texttt{\$\{bmc\_username\}} & {\small \# BMC username for use by IPMI} \\
& \texttt{\$\{bmc\_password\}} & {\small \# BMC password for use by IPMI} \\
& \texttt{\$\{c\_ip[0]\}}, \texttt{\$\{c\_ip[1]\}}, ... & {\small \# Desired compute node addresses} \\
& \texttt{\$\{c\_bmc[0]\}}, \texttt{\$\{c\_bmc[1]\}}, ... & {\small \# BMC addresses for computes} \\
& \texttt{\$\{c\_mac[0]\}}, \texttt{\$\{c\_mac[1]\}}, ... & {\small \# MAC addresses for computes} \\
& \texttt{\$\{compute\_regex\}} & {\small \# Regex for matching compute node names (e.g. c*)} \\
\end{tabular}

\vspace*{0.2cm}
\noindent {Optional:} 
\vspace*{0.1cm}

\begin{tabular}{@{}>{\textbullet}l p{7cm} l}
& \texttt{\$\{mgs\_fs\_name\}} & {\small \# Lustre MGS mount name} \\
& \texttt{\$\{sms\_ipoib\}} & {\small \# IPoIB address for SMS server} \\
& \texttt{\$\{ipoib\_netmask\}} & {\small \# Subnet netmask for internal IPoIB} \\
& \texttt{\$\{c\_ipoib[0]\}}, \texttt{\$\{c\_ipoib[1]\}}, ... & {\small \# IPoIB addresses for computes} \\
\end{tabular}




% Base Operating System --------------------------------------------

\section{Install Base Operating System (BOS)}
In an external setting, installing the desired BOS on a {\em master} host would
typically involve booting from a DVD iso image on a new server.  However, on
the Zeus cluster, master nodes can be network booted with a minimal OS install
using an external {\em Warewulf} instance. The Warewulf configuration
provisions the OS into a ramdisk, configures local package repositories to
point to the FSP OBS server and underlying distro repos, and mounts an
NFS \$HOME file system. With this setup, all that is required to achieve a
clean BOS intall on a master host is to perform a host reboot. If, however, you
are installing a new system, insert the \baseOS{} DVD, power cycle the host,
and follow the distro provided directions to install the BOS on your chosen
{\em master} host.  


% ------------------------------------------------------------------

\section{Install \FSP{} Components} \label{sec:basic_install}
With the BOS installed and booted, the next step is to add desired FSP packages
onto the {\em master} server in order to provide provisioning and resource
management services for the rest of the cluster. The following subsections
highlight this process.


\subsection{Enable \FSP{} repository for local use} \label{sec:enable_repo}
To begin, enable use of the FSP repository by adding it to the local list of
available package repositories.  Note that this requires network access from
your {\em master} server to the FSP repository, or alternatively, that the FSP
repository be mirrored locally. The example below illustrates the addition
using the internal FSP checkpoint repository for the \baseOS{} distribution.


% begin_fsp_run
% fsp_validation_comment Enable FSP Repo(s)
\begin{lstlisting}[language=bash,keywords={}]
[master]$ export FSP_mirror=http://tcgsw-obs.pdx.intel.com:82
[master]$ wget -P /etc/yum.repos.d $FSP_mirror/ForestPeak:/15.27:/Factory/CentOS-7.1_Intel/ForestPeak:15.27:Factory.repo
\end{lstlisting}
% end_fsp_run

% Dev testing for Factory
% [master]$ wget -P /etc/yum.repos.d $FSP_mirror/ForestPeak:/15.16/CentOS-7.1_Intel/ForestPeak:15.16.repo

In addition to the \FSP{} package repository, the {\em master} host also
requires access to the standard base OS distro repositories in order to resolve
necessary dependencies. For \baseOS{}, the requirements are to have access to
both the base OS and EPEL repositories (mirrors freely available online):

\begin{itemize*}
\item CentOS-7 - Base 7.1.1503
  (e.g. \href{http://mirror.centos.org/centos/7.1.1503/os/x86\_64}
             {\color{blue}{http://mirror.centos.org/centos/7.1.1503/os/x86\_64}} )
%\item EPEL 7 (e.g. {\color{blue}
%  \url{http://download.fedoraproject.org/pub/epel/7}} )
\item EPEL 7 (e.g. \href{http://download.fedoraproject.org/pub/epel/7/x86\_64}
                        {\color{blue}{http://download.fedoraproject.org/pub/epel/7/x86\_64}} )
\end{itemize*}

% begin_fsp_run
% % fsp_validation_comment - format space for /opt/fsp (CI only command)
% [master]$ mkfs -F -t ext4 /dev/sda1
% [master]$ mkdir -p /opt/fsp
% [master]$ mount /dev/sda1 /opt/fsp
% end_fsp_run  

\subsection{Add provisioning services on {\em master} node}
With the FSP repo enabled, we can now begin adding desired components onto the
{\em master} server. The FSP repo provides a number of aliases that group
logical components togetther in order to help aid in this process. To add
support for provisioning services, the following commands illustrate addition
of a common base package followed by the Warewulf provisioning system.


% begin_fsp_run 
% fsp_validation_comment - allow sudo without tty (CI only)
% [master]$ perl -pi -e "s/^Defaults\s+requiretty/# Defaults requiretty/" /etc/sudoers 
% end_fsp_run

\subsection{Add resource management services on {\em master} node} \label{sec:add_rm}
\OHPC{} provides multiple options for distributed resource management. 
The following command adds the \SLURM{} workload manager server components to the
chosen {\em master} host. Note that client-side components will be added to
the corresponding compute image in a subsequent step.

% begin_ohpc_run
% ohpc_comment_header Add resource management services on master node \ref{sec:add_rm}
\begin{lstlisting}[language=bash,keywords={}]
# Install slurm server meta-package
[sms](*\#*) (*\install*) ohpc-slurm-server

# Identify resource manager hostname on master host
[sms](*\#*) perl -pi -e "s/ControlMachine=\S+/ControlMachine=${sms_name}/" /etc/slurm/slurm.conf
\end{lstlisting}
% end_ohpc_run

\begin{center}
\begin{tcolorbox}[]
  \small SLURM requires enumeration of the physical hardware characteristics
  for compute nodes under its control. In particular, three configuration
  parameters combine to define consumable compute resources: {\em Sockets},
  {\em CoresPerSocket}, and {\em ThreadsPerCore}. The default configuration
  file provided via \OHPC{} assumes dual-socket, 8 cores per socket, and two
  threads per core for this 4-node example. If this does not reflect your local
  hardware, please update the configuration file at
  \path{/etc/slurm/slurm.conf} accordingly to match your particular hardware.
\end{tcolorbox}
\end{center}

% begin_ohpc_run
% ohpc_comment_header Update node configuration for slurm.conf
% ohpc_command if [[ ${update_slurm_nodeconfig} -eq 1 ]];then
% ohpc_indent 5
% ohpc_command perl -pi -e "s/^NodeName=.+$/#/" /etc/slurm/slurm.conf
% ohpc_command perl -pi -e "s/^NodeName=c\S+/NodeName=c[1-$num_computes]/" /etc/slurm/slurm.conf
% ohpc_command echo -e ${slurm_node_config} >> /etc/slurm/slurm.conf
% ohpc_indent 0
% ohpc_command fi
% end_ohpc_run

Other versions of this guide are available that describe installation of alternate
resource management systems, and they can be found in the \texttt{docs-ohpc}
package.



\subsection{Add \InfiniBand{} support services on {\em master} node} \label{sec:add_ofed}

The following command adds OFED support using distro-provided drivers to the
chosen {\em master} host. 

% begin_fsp_run
% fsp_validation_comment Add ofed
\begin{lstlisting}[language=bash,keywords={}]
[master]$ (*\groupinstall*) "InfiniBand Support"

# Load IB drivers
[master]$ service rdma start
\end{lstlisting}
% end_fsp_run

With the \InfiniBand{} drivers included, you can also enable IPoIB functionality
which provides a mechanism to send IP packets over the IB network. If you plan
to mount a \Lustre{} file system over \InfiniBand{} (see \S\ref{sec:lustre_client}
for additional details), then having IPoIB enabled is a requirement for the
\Lustre{} client. \FSP{} provides a template configuration file to aid in setting up
an {\em ib0} interface on the {\em master} host. To use, copy the template
provided and update the \texttt{master\_ipoib} and
\texttt{ipoib\_netmask} entries to match local desired settings (alter ib0
naming as appropriate if system contains dual-ported or multiple HCAs). 

% begin_fsp_run
% fsp_validation_comment Enable ib0
\begin{lstlisting}[language=bash,literate={-}{-}1,keywords={},upquote=true]
[master]$ cp /opt/fsp/pub/examples/network/centos/ifcfg-ib0 /etc/sysconfig/network-scripts

# Define local IPoIB address and netmask
[master]$ perl -pi -e 's/master_ipoib/<master_ipoib>/' /etc/sysconfig/network-scripts/ifcfg-ib0
[master]$ perl -pi -e 's/ipoib_netmask/<ipoib_netmask>/' /etc/sysconfig/network-scripts/ifcfg-ib0

# Initiate ib0
[master]$ ifup ib0
\end{lstlisting}
% end_fsp_run

\subsection{Add Cluster Checker} \label{sec:add_clck}
The \Intel{} Cluster Checker provides a convenient suite of diagnostics that
can be used to aid in isolating hardware and software problems on an installed
cluster. This package can be optionally added using the command below. Note
that a valid license file will also need to be installed separately in order to
use this software.

% begin_fsp_run
% fsp_validation_comment Add CLCK
\begin{lstlisting}[language=bash,keywords={}]
[master]$ (*\install*) intel-clck-fsp
\end{lstlisting}
% end_fsp_run


\subsection{Complete basic Warewulf setup for {\em master} node}

At this point, all of the packages necessary to use \Warewulf{} on the {\em
  master} host should be installed.  Next, we need to update several
configuration files in order to allow \Warewulf{} to work with \baseOS{} and to support
local provisioning using the second ethernet interface (assumed to be {\em
  eno2} in this example).  Specific steps are as
follows:

% begin_fsp_run
% fsp_validation_comment Complete basic Warewulf setup for master node
%\begin{verbatim}

\begin{lstlisting}[language=bash,literate={-}{-}1,keywords={},upquote=true,keepspaces]
# Configure Warewulf provisioning to use eno2 interface
[master]$ perl -pi -e 's/^network device = eth1/network device = eno2/' /etc/warewulf/provision.conf

# Enable tftp service for compute node image distribution
[master]$ perl -pi -e "s/^\s+disable\s+= yes/ disable = no/" /etc/xinetd.d/tftp

# Enable http access for Warewulf cgi-bin directory to support newer apache syntax
[master]$ export MODFILE=/etc/httpd/conf.d/warewulf-httpd.conf
[master]$ perl -pi -e "s/cgi-bin>\$/cgi-bin>\n Require all granted/" $MODFILE

[master]$ perl -pi -e "s/Allow from all/Require all granted/" $MODFILE
[master]$ perl -ni -e "print unless /^\s+Order allow,deny/" $MODFILE

# Define eno2 interface for provisioning
[master]$ ifconfig eno2 <master_ip> netmask 255.255.255.0 up

# Restart relevant services to support provisioning
[master]$ systemctl restart xinetd  
[master]$ systemctl restart mariadb 
[master]$ systemctl restart httpd   
[master]$ systemctl restart dhcpd
\end{lstlisting}
%\end{verbatim}
% end_fsp_run

% begin_fsp_run 
% fsp_validation_comment - time update
% [master]$ TIME_T2=$(date +%s) ; TIME_DELTA=$(($TIME_T2-$TIME_T1)) ; TIME_T1=$TIME_T2
% [master]$ echo " "; echo "END TIMING = $TIME_DELTA (secs)"
% end_fsp_run


\subsection{Define {\em compute} image for provisioning}

With the provisioning services enabled, the next step is to define and
customize a system image that can subsequently be used to provision one or more
{\em compute} nodes. The following subsections highlight this process.

\subsubsection{Build initial BOS image}

The \FSP{} build of \Warewulf{} includes specific enhancements enabling support for
\baseOS{}. The following steps illustrate the process to build a minimal, default
image for use with \Warewulf{}.  We begin by creating a directory structure on the 
{\em master} that will represent the root filesystem of the compute node. The 
default location for this in \FSP{} is in
\texttt{/opt/fsp/admin/images/centos7.1}. Note that \Warewulf{} assumes access to
an external repository (mirror.centos.org) during the \texttt{wwmkchroot}
process which requires external network connectivity on your chosen {\em
  master} host. If this is not available, or if you prefer to access a locally cached
mirror, the steps below include an example to update the OS mirror used by
\texttt{wwmkchroot} to use an alternate location defined by a
\texttt{BOS\_MIRROR} environment variable.

% begin_fsp_run
% fsp_validation_comment Create compute image for Warewulf
\begin{lstlisting}[language=bash,literate={-}{-}1,keywords={},upquote=true,keepspaces]
# Define chroot location 
[master]$ export CHROOT=/opt/fsp/admin/images/centos7.1

# Override default OS repository (optional) - set BOS_MIRROR variable to desired location
[master]$ perl -pi -e "s#^YUM_MIRROR=(\S+)#YUM_MIRROR=$BOS_MIRROR#" \
     /usr/libexec/warewulf/wwmkchroot/centos-7.tmpl

# Build initial chroot image
[master]$ wwmkchroot centos-7 $CHROOT
\end{lstlisting}
% end_fsp_run

\subsubsection{Add \FSP{} components}

The \texttt{wwmkchroot} process used in the previous step is designed to
provide a minimal \baseOS{} configuration. Next, we add additional components to
include resource management client services, \InfiniBand{} drivers, and other
additional packages to support the default \FSP{} environment.  This process uses
the \texttt{chroot} command to augment the base provisioning image and will
access the BOS and \FSP{} repositories to resolve package install requests. To
access the remote repositories by hostname (and not IP addresses), the chroot
environment needs to be updated to enable DNS resolution. Assuming that
the {\em master} host has a working DNS configuration in place, the chroot environment can
be updated with a copy of the configuration as follows:

% begin_fsp_run
\begin{lstlisting}[language=bash,literate={-}{-}1,keywords={},upquote=true]
[master]$ cp -p /etc/resolv.conf $CHROOT/etc/resolv.conf
\end{lstlisting}
% end_fsp_run

Prior to using the \FSP{} repository to add desired yum groups, it is
necessary to fix the cached \texttt{repomod.xml} file within the chroot
environment (similar to the initial modification made for the {\em master} host
in \S\ref{sec:enable_repo}). This can be accomplished as follows:

% begin_fsp_run
% fsp_validation_comment - fix grouplist for chroot
\begin{lstlisting}[language=bash,literate={-}{-}1,keywords={},upquote=true]
[master]$ chroot $CHROOT yum grouplist
[master]$ perl -i -pe 's/type="group"/type="group_gz"/' $CHROOT/var/cache/yum/x86_64/7/ForestPeak_*/repomd.xml
\end{lstlisting}
% end_fsp_run

Now, we can include additional components to the compute instance using
\texttt{yum} to install into the chroot location defined previously:

% begin_fsp_run
% fsp_validation_comment Add FSP components to compute instance
\begin{lstlisting}[language=bash,literate={-}{-}1,keywords={},upquote=true]
# Add SLURM client support
[master]$ (*\groupchrootinstall*) fsp-slurm-client

# Add IB support and enable
[master]$ (*\groupchrootinstall*) "InfiniBand Support"
[master]$ chroot $CHROOT systemctl enable rdma

# Add Network Time Protocol (NTP) support
[master]$ (*\chrootinstall*) ntp

# Include FSP default modules user environment
[master]$ (*\chrootinstall*) lmod-defaults-intel-fsp
\end{lstlisting}
% end_fsp_run


% begin_fsp_run 
% fsp_validation_comment - Restart openibd (CI only command)
% [master]$ service openibd restart
% end_fsp_run

\subsubsection{Customize system configuration} \label{sec:master_customization}

Prior to assembling the image, it is advantageous to perform any additional
customizations within the chroot environment created for the desired {\em
  compute} instance. The following steps document the process to add a local
{\em ssh} key created by \Warewulf{} to support remote access, identify the
resource manager server, configure NTP for compute resources, and enable \NFS{}
mounting of a \$HOME file system and the public \FSP{} install path
(\texttt{/opt/fsp/pub}) that will be hosted by the {\em master} host in this
example configuration.  The \NFS{} exporting options use an address/netmask
combination to limit the export scope to the defined compute nodes.

% begin_fsp_run
% fsp_validation_comment Add ssh key and enable NFS mount of /opt/fsp/pub
\begin{lstlisting}[language=bash,literate={-}{-}1,keywords={},upquote=true]
# add new cluster key to base image
[master]$ cat ~/.ssh/cluster.pub >> $CHROOT/root/.ssh/authorized_keys

# add NFS client mounts of /home and /opt/fsp/pub to base image
[master]$ echo "<master_ip>:/home /home nfs nfsvers=3,rsize=1024,wsize=1024,cto 0 0" >> $CHROOT/etc/fstab
[master]$ echo "<master_ip>:/opt/fsp/pub /opt/fsp/pub nfs nfsvers=3 0 0" >> $CHROOT/etc/fstab

# Identify resource manager server IP in SLURM config file on computes
[master]$ perl -pi -e "s/ControlMachine=\S+/ControlMachine=<master_ip>/" $CHROOT/etc/slurm/slurm.conf

# Identify resource manager hostname on master host
[master]$ perl -pi -e "s/ControlMachine=\S+/ControlMachine=<master_hostname>/" /etc/slurm/slurm.conf

# Export /home and FSP public packages from master server to cluster compute nodes
[master]$ echo "/home <internal_netmask>(rw,no_subtree_check,fsid=10,no_root_squash)" >> /etc/exports
[master]$ echo "/opt/fsp/pub <internal_netmask>(ro,no_subtree_check,fsid=11)" >> /etc/exports
[master]$ exportfs -a
[master]$ systemctl restart nfs

# Enable NTP time service on computes and identify master host as local NTP server
[master]$ chroot $CHROOT systemctl enable ntpd
[master]$ echo "server <master_ip>" >> $CHROOT/etc/ntp.conf
\end{lstlisting}
% end_fsp_run

\subsubsection{Additional Customizations ({\em optional})}

This section highlights common additional customizations that
can {\em optionally} be applied to the
local cluster environment. These customizations include:

\begin{itemize*}
\item Increase memlimits to support \InfiniBand{}
\item Restrict ssh access to compute resources
\item Add \Lustre{} client
\end{itemize*}

\noindent Details on the steps required for each of these customizations are
discussed further in the following sections.

\paragraph{Increase locked memory limits}

In order to utilize \InfiniBand{} or Omni-Path as the underlying high speed interconnect, it is
generally necessary to increase the locked memory settings for system
users. This can be accomplished by updating
the \texttt{/etc/security/limits.conf} file and this should be performed within
the {{\em compute} image and on all job submission hosts. In this recipe, jobs
are submitted from the {\em master} host, and the following commands can be
used to update the maximum locked memory settings on both the master host and
the compute image:

% begin_ohpc_run
% ohpc_validation_newline 
% ohpc_validation_comment Update memlock settings
\begin{lstlisting}[language=bash,keywords={},upquote=true]
# Update memlock settings on master
[sms](*\#*) perl -pi -e 's/# End of file/\* soft memlock unlimited\n$&/s' /etc/security/limits.conf
[sms](*\#*) perl -pi -e 's/# End of file/\* hard memlock unlimited\n$&/s' /etc/security/limits.conf

# Update memlock settings within compute image
[sms](*\#*) perl -pi -e 's/# End of file/\* soft memlock unlimited\n$&/s' $CHROOT/etc/security/limits.conf
[sms](*\#*) perl -pi -e 's/# End of file/\* hard memlock unlimited\n$&/s' $CHROOT/etc/security/limits.conf
\end{lstlisting}
% end_ohpc_run


 

% begin_fsp_run
% fsp_validation_comment Update memlock on master
\begin{lstlisting}[language=bash,keywords={},upquote=true]
# Update memlock settings on master
[master]$ echo "* soft memlock unlimited" >> /etc/security/limits.conf
[master]$ echo "* hard memlock unlimited" >> /etc/security/limits.conf

# Update memlock settings within compute image
[master]$ echo "* soft memlock unlimited" >> $CHROOT/etc/security/limits.conf
[master]$ echo "* hard memlock unlimited" >> $CHROOT/etc/security/limits.conf
\end{lstlisting}
% end_fsp_run

%# Enable larger locked memory limits for SLURM daemon
%[master]$ echo "ulimit -l unlimited" >> $CHROOT/etc/sysconfig/slurm

\paragraph{Enable ssh control via resource manager} 

An additional optional customization that is recommended is to
restrict \texttt{ssh} access on compute nodes to only allow access by users who
have an active job associated with the node. This can be enabled via the use of
a pluggable authentication module (PAM) provided as part of the SLURM package
installs. To enable this feature within the {\em compute} image, issue the following:




% begin_fsp_run
% fsp_validation_comment Enable slurm pam module
\begin{lstlisting}[language=bash,keywords={},upquote=true]
[master]$ echo "account    required     pam_slurm.so" >> $CHROOT/etc/pam.d/sshd
\end{lstlisting}
% end_fsp_run

\paragraph{Add \Lustre{} client} \label{sec:lustre_client}

To add \Lustre{} client support on the cluster, it necessary to install the client
and associated modules on each host needing to access a \Lustre{} file system. In
this recipe, it is assumed that the \Lustre{} file system is hosted by servers
that are pre-existing and are not part of the install process. Outlining the
variety of \Lustre{} client mounting options is beyond the scope of this document,
%(please consult \Lustre{} documentation for more details on failover configuration
%support and networking options), 
but the general requirement is to add a mount entry for the desired file system
that defines the management server (MGS) and underlying network transport
protocol. To add client mounts on both the {\em master} server and {\em
compute} image, the following commands can be used. Note that the \Lustre{} file
system to be mounted is identified by the \texttt{\$\{mgs\_fs\_name\}} variable. 
In this example, the file system is configured to be mounted locally
as \path{/mnt/lustre}.

%% % SLES-specific addition
%% Additionally, note that your default \baseOS{} environment may not allow loading of
%% the necessary \Lustre{} kernel modules. Consequently, the example below includes
%% steps which update the {\texttt /etc/modprobe.d/10-unsupported-modules.conf}
%% file to allow loading of the necessary modules.

% begin_fsp_run
% fsp_validation_comment Install Lustre client on master
\begin{lstlisting}[language=bash,keywords={},upquote=true]
# Add Lustre client software to master host
[master]$ (*\install*) lustre-client-fsp lustre-client-fsp-modules
\end{lstlisting}
% end_fsp_run

% begin_fsp_run
% fsp_validation_comment Enable lustre in WW compute image
\begin{lstlisting}[language=bash,keywords={},upquote=true]
# Include Lustre client software in compute image
[master]$ (*\chrootinstall*) lustre-client-fsp lustre-client-fsp-modules

# Include mount point and file system mount in compute image
[master]$ mkdir $CHROOT/mnt/lustre
[master]$ echo "<mgs_fs_name> /mnt/lustre lustre defaults,_netdev,localflock 0 0" >> $CHROOT/etc/fstab
\end{lstlisting}
% end_fsp_run

The default underlying network type used by \Lustre{} is {\em tcp}. If your
external \Lustre{} file system is to be mounted using a network type other than
tcp, additional configuration files are necessary to identify the desired
network type. The example below illustrates creation of modprobe config files
instructing \Lustre{} to use an \InfiniBand{} network with the {\em o2ib} LNET driver
attached to \texttt{ib0}. Note that these modifications are made to both the
{\em master} host and {\em compute} image.

% begin_fsp_run
% fsp_validation_comment Enable o2ib for Lustre
\begin{lstlisting}[language=bash,keywords={},upquote=true]
[master]$ echo "options lnet networks=o2ib(ib0)" >> /etc/modprobe.d/lustre.conf
[master]$ echo "options lnet networks=o2ib(ib0)" >> $CHROOT/etc/modprobe.d/lustre.conf
\end{lstlisting}
% end_fsp_run

With the \Lustre{} config complete, the client can be mounted on the {\em master}
host as follows:
% begin_fsp_run
% fsp_validation_comment mount Lustre client on master
\begin{lstlisting}[language=bash,keywords={},upquote=true]
[master]$ mkdir /mnt/lustre
[master]$ mount -t lustre -o localflock <mgs_fs_name> /mnt/lustre
\end{lstlisting}
% end_fsp_run

%% [FSP-469]: disabling ORCM

%% [FSP-469] \paragraph{Enable ORCM RAS subsystem} 
%% [FSP-469] \OHPC{} includes a monitoring capability that can optionally be installed to
support unified data collection across the cluster. This monitoring system is
based on \ORCM{}, the open resilient cluster manager project, and provides a
convenient way to aggregate key node metrics including environmental sensors,
processor power, and resource utilization into a centralized database for
historical logging and analysis. In a typical HPC configuration, \ORCM{} data
collection relies on having IPMI access to the baseboard management controller
(BMC) on each back-end compute resource. The data collection itself can either
be performed using {\em in-band} methods where a daemon is instantiated on each
compute resource or via {\em out-of-band} methods using external IPMI access on
a management network. To include all the components necessary to enable an \ORCM{}
server running on the {\em master} host, issue the following:

% begin_ohpc_run
% ohpc_validation_comment Install ORCM server
\begin{lstlisting}[language=bash]
[sms](*\#*) (*\groupinstall*) ohpc-orcm-server
\end{lstlisting}
% end_ohpc_run

In order to coalesce sensor data into a central repository, \ORCM{} requires
data {\em aggregators} that are responsible for collecting data from one or
more monitored resources. For a small cluster, a common configuration is to
have a single aggregator defined and, in this example, we choose the {\em master} host
as the data aggregator. In order to store historical logging data, \ORCM{}
includes support to store aggregated data into a centralized database. In this
example, we will create an \ORCM{} database using {\em postgres} on the chosen
{\em master} host. The necessary configuration file changes and commands to create the
underlying database for \ORCM{} are highlighted below. Note that a generic DB
password of ``orcmpassword'' is used; this should be altered locally to a
site-specific credential. If the DB password is changed, a corresponding update
for the new credential should also be placed in the final
\texttt{/etc/sysconfig/orcmd} file.

%% [FSP-469] 
%% [FSP-469] % begin_fsp_run
%% [FSP-469] % fsp_validation_comment Configure ORCM aggregator and DB
%% [FSP-469] \begin{lstlisting}[language=bash,keywords={},upquote=true,keepspaces]
%% [FSP-469] # Set master host as data aggregator
%% [FSP-469] [master]$ perl -pi -e "s/aggregator_hostname/localhost/" /opt/open-rcm/etc/orcm-site.xml
%% [FSP-469] [master]$ perl -pi -e "s/Servername = (<.+>)/Servername = <master_ip>/" /etc/orcmdb_psql.ini
%% [FSP-469] 
%% [FSP-469] # Update BMC credentials to match local values
%% [FSP-469] 
%% [FSP-469] [master]$ perl -pi -e "s/your-bmc-username/<bmc_username>/" /opt/open-rcm/etc/openmpi-mca-params.conf
%% [FSP-469] [master]$ perl -pi -e "s/your-bmc-password/<bmc_password>/" /opt/open-rcm/etc/openmpi-mca-params.conf
%% [FSP-469] 
%% [FSP-469] # Initialize postgres DB driver
%% [FSP-469] [master]$ odbcinst -i -d -f /etc/psql_odbc_driver.ini
%% [FSP-469] [master]$ odbcinst -i -s -h -f /etc/orcmdb_psql.ini 
%% [FSP-469] 
%% [FSP-469] # Initialize postgres
%% [FSP-469] [master]$ postgresql-setup initdb
%% [FSP-469] 
%% [FSP-469] # Update postgres config files using ORCM templates
%% [FSP-469] [master]$ cp /etc/postgresql.orcm.conf /var/lib/pgsql/data/postgresql.conf
%% [FSP-469] [master]$ cp /etc/pg_hba.orcm.conf /var/lib/pgsql/data/pg_hba.conf
%% [FSP-469] 
%% [FSP-469] # Restart postgres
%% [FSP-469] [master]$ service postgresql start
%% [FSP-469] 
%% [FSP-469] # Create ORCM DB
%% [FSP-469] [master]$ sudo -u postgres psql --command "CREATE USER orcmuser WITH SUPERUSER PASSWORD 'orcmpassword';"
%% [FSP-469] [master]$ sudo -u postgres createdb -O orcmuser orcmdb
%% [FSP-469] [master]$ sudo -u postgres psql --username=orcmuser --dbname=orcmdb -f /etc/orcmdb_psql.sql
%% [FSP-469] 
%% [FSP-469] # Start ORCM on master host (requires both orcmsched and orcmd)
%% [FSP-469] [master]$ systemctl start orcmsched
%% [FSP-469] [master]$ systemctl start orcmd
%% [FSP-469] \end{lstlisting}
%% [FSP-469] %% [FSP-469] % end_fsp_run
%% [FSP-469] 
%% [FSP-469] To use the {\em in-band} method of collecting data from compute resources, the
\ORCM{} client daemon is required to be available on each compute node. This can
be added to the compute image as follows:

% begin_fsp_run
% fsp_validation_comment Install ORCM client on compute image
\begin{lstlisting}[language=bash,keywords={},upquote=true]
[master]$ (*\groupchrootinstall*) fsp-orcm-client

# Set master host as data aggregator
[master]$ perl -pi -e "s/aggregator_hostname/<master_ip>/" $CHROOT/opt/open-rcm/etc/orcm-site.xml

# Enable orcmd client daemon on boot
[master]$ chroot $CHROOT systemctl enable orcmd
\end{lstlisting}
% end_fsp_run


\subsubsection{Import files} \label{sec:file_import}

The \Warewulf{} system includes functionality to import arbitrary files from the
provisioning server for distribution to managed hosts. This is one way
to distribute user credentials across to {\em compute} hosts. To
import local file-based credentials, issue the following:

% begin_fsp_run
% fsp_validation_comment Import credentials
\begin{lstlisting}[language=bash,literate={-}{-}1,keywords={},upquote=true]
[master]$ wwsh file import /etc/passwd                                                                                                       
[master]$ wwsh file import /etc/group
[master]$ wwsh file import /etc/shadow 
\end{lstlisting}
% \end_fsp_run

Similarly, to import the cryptographic key that is required by the {\em munge}
authentication library to be available on every host in the resource management
pool, issue the following:

% begin_fsp_run
\begin{lstlisting}[language=bash,literate={-}{-}1,keywords={},upquote=true]
[master]$ wwsh file import /etc/munge/munge.key
\end{lstlisting}
% \end_fsp_run

Finally, to add support for controlling IPoIB interfaces, \FSP{} includes a
template file for \Warewulf{} that can be imported and used later to provision
\texttt{ib0} network settings.

% begin_fsp_run
\begin{lstlisting}[language=bash,literate={-}{-}1,keywords={},upquote=true]
[master]$ wwsh file import /opt/fsp/pub/examples/network/centos/ifcfg-ib0.ww
[master]$ wwsh -y file set ifcfg-ib0.ww --path=/etc/sysconfig/network-scripts/ifcfg-ib0
\end{lstlisting}
% \end_fsp_run

% begin_fsp_run 
% fsp_validation_comment - time update
% [master]$ TIME_T2=$(date +%s) ; TIME_DELTA=$(($TIME_T2-$TIME_T1)) ; TIME_T1=$TIME_T2
% [master]$ echo " "; echo "END TIMING = $TIME_DELTA (secs)"
% end_fsp_run

\subsection{Assemble bootstrap image}

\Warewulf{} requires a bootstrap image that is used to initialize the provisioning
process. This bootstrap image includes the runtime kernel and associated
modules. The following commands highlight the inclusion of additional drivers
and creation of the bootstrap image based on the running kernel.

% begin_fsp_run
% fsp_validation_comment Create bootstrap image for Warewulf
\begin{lstlisting}[language=bash,literate={-}{-}1,keywords={},upquote=true]
# (Optional) Include Lustre drivers; needed if enabling Lustre client on computes
[master]$ export WW_CONF=/etc/warewulf/bootstrap.conf
[master]$ echo "drivers += updates/kernel/" >> $WW_CONF

# Build bootstrap image
[master]$ wwbootstrap `uname -r`
\end{lstlisting}
% end_fsp_run

\subsubsection{Assemble Virtual Node File Sytem (VNFS) image}

With the local site customizations in place, the following step uses the
\texttt{wwvnfs} command to assemble a VNFS capsule from the chroot environment
defined for the {\em compute} instance. 

% begin_fsp_run
% fsp_validation_comment Assemble VNFS
\begin{lstlisting}[language=bash,literate={-}{-}1,keywords={},upquote=true]
[master]$ wwvnfs -y --chroot $CHROOT
\end{lstlisting}
% end_fsp_run

\subsubsection{Register nodes for provisioning}

In preparation for provisioning, we can now define the desired network settings
for four example compute nodes with the underlying provisioning system and
restart the \texttt{dhcp} service. Note the use of variable names for the
desired compute node IP and MAC addresses which should be modified to
accommodate local settings and hardware. Included in these steps are commands
to enable \Warewulf{} to manage IPoIB settings and corresponding definitions of
IPoIB addresses for the compute nodes. This is typically optional unless you
are planning to include a \Lustre{} client mount over \InfiniBand{}.  The final step
in this process associates the VNFS image assembled in previous steps with the
newly defined compute nodes, utilizing the user credential files and munge key
that were imported in \S\ref{sec:file_import}.

% begin_fsp_run
% fsp_validation_comment Add hosts to cluster

\begin{lstlisting}[language=bash,keywords={},upquote=true,basicstyle=\footnotesize\ttfamily]
# Define four compute nodes and network settings 
[master]$ wwsh -y node new c1 --ipaddr=<c1_ip> --hwaddr=<c1_mac> 
[master]$ wwsh -y node new c2 --ipaddr=<c2_ip> --hwaddr=<c2_mac> 
[master]$ wwsh -y node new c3 --ipaddr=<c3_ip> --hwaddr=<c3_mac> 
[master]$ wwsh -y node new c4 --ipaddr=<c4_ip> --hwaddr=<c4_mac> 

# Define provisioning image for hosts
[master]$ wwsh -y provision set c[1-4] --vnfs=centos7.1 --bootstrap=`uname -r` \
     --files=dynamic_hosts,passwd,group,shadow,munge.key 

# Optionally define IPoIB network settings (required if planning to mount Lustre* over IB)
[master]$ wwsh -y node set c1 -D ib0 --ipaddr=<c1_ipoib> --netmask=<ipoib_netmask>
[master]$ wwsh -y node set c2 -D ib0 --ipaddr=<c2_ipoib> --netmask=<ipoib_netmask>
[master]$ wwsh -y node set c3 -D ib0 --ipaddr=<c3_ipoib> --netmask=<ipoib_netmask>
[master]$ wwsh -y node set c4 -D ib0 --ipaddr=<c4_ipoib> --netmask=<ipoib_netmask>
[master]$ wwsh -y provision set c[1-4] --fileadd=ifcfg-ib0.ww

# Restart dhcp / update PXE
[master]$ service dhcpd restart
[master]$ wwsh pxe update
\end{lstlisting}
% end_fsp_run

% begin_fsp_run 
% fsp_validation_comment - optionally add 4 more hosts
% [master]$ if [[ $USE_DYNAMIC_HOSTNUM -eq 1 && $NUM_COMPUTES -ge 8 ]];then
% [master]$    wwsh -y node new c5 --ipaddr=<c5_ip> --hwaddr=<c5_mac> 
% [master]$    wwsh -y node new c6 --ipaddr=<c6_ip> --hwaddr=<c6_mac> 
% [master]$    wwsh -y node new c7 --ipaddr=<c7_ip> --hwaddr=<c7_mac> 
% [master]$    wwsh -y node new c8 --ipaddr=<c8_ip> --hwaddr=<c8_mac> 
% [master]$    wwsh -y node set c5 -D ib0 --ipaddr=<c5_ipoib> --netmask=<ipoib_netmask>
% [master]$    wwsh -y node set c6 -D ib0 --ipaddr=<c6_ipoib> --netmask=<ipoib_netmask>
% [master]$    wwsh -y node set c7 -D ib0 --ipaddr=<c7_ipoib> --netmask=<ipoib_netmask>
% [master]$    wwsh -y node set c8 -D ib0 --ipaddr=<c8_ipoib> --netmask=<ipoib_netmask>
% [master]$    wwsh -y provision set c[5-8] --vnfs=sles12 --bootstrap=`uname -r` --files=dynamic_hosts,passwd,group,shadow,munge.key 
% [master]$    wwsh -y provision set c[5-8] --fileadd=ifcfg-ib0.ww     
% [master]$    service dhcpd restart
% [master]$    wwsh pxe update
% [master]$ fi
% end_fsp_run

% begin_fsp_run 
% fsp_validation_comment - optionally add 8 more hosts
% [master]$ if [[ $USE_DYNAMIC_HOSTNUM -eq 1 && $NUM_COMPUTES -ge 16 ]];then
% [master]$    wwsh -y node new c9  --ipaddr=<c9_ip>  --hwaddr=<c9_mac> 
% [master]$    wwsh -y node new c10 --ipaddr=<c10_ip> --hwaddr=<c10_mac> 
% [master]$    wwsh -y node new c11 --ipaddr=<c11_ip> --hwaddr=<c11_mac> 
% [master]$    wwsh -y node new c12 --ipaddr=<c12_ip> --hwaddr=<c12_mac> 

% [master]$    wwsh -y node new c13 --ipaddr=<c13_ip> --hwaddr=<c13_mac> 
% [master]$    wwsh -y node new c14 --ipaddr=<c14_ip> --hwaddr=<c14_mac> 
% [master]$    wwsh -y node new c15 --ipaddr=<c15_ip> --hwaddr=<c15_mac> 
% [master]$    wwsh -y node new c16 --ipaddr=<c16_ip> --hwaddr=<c16_mac> 

% [master]$    wwsh -y node set c9  -D ib0 --ipaddr=<c9_ipoib>  --netmask=<ipoib_netmask>
% [master]$    wwsh -y node set c10 -D ib0 --ipaddr=<c10_ipoib> --netmask=<ipoib_netmask>
% [master]$    wwsh -y node set c11 -D ib0 --ipaddr=<c11_ipoib> --netmask=<ipoib_netmask>
% [master]$    wwsh -y node set c12 -D ib0 --ipaddr=<c12_ipoib> --netmask=<ipoib_netmask>
% [master]$    wwsh -y node set c13 -D ib0 --ipaddr=<c13_ipoib> --netmask=<ipoib_netmask>
% [master]$    wwsh -y node set c14 -D ib0 --ipaddr=<c14_ipoib> --netmask=<ipoib_netmask>
% [master]$    wwsh -y node set c15 -D ib0 --ipaddr=<c15_ipoib> --netmask=<ipoib_netmask>
% [master]$    wwsh -y node set c16 -D ib0 --ipaddr=<c16_ipoib> --netmask=<ipoib_netmask>

% [master]$    wwsh -y provision set c* --vnfs=centos-7 --bootstrap=`uname -r` --files=dynamic_hosts,passwd,group,shadow,munge.key 
% [master]$    wwsh -y provision set c* --fileadd=ifcfg-ib0.ww     
% [master]$    service dhcpd restart
% [master]$    wwsh pxe update
% [master]$ fi
% end_fsp_run

% begin_fsp_run 
% fsp_validation_comment - time update
% [master]$ TIME_T2=$(date +%s) ; TIME_DELTA=$(($TIME_T2-$TIME_T1)) ; TIME_T1=$TIME_T2
% [master]$ echo " "; echo "END TIMING = $TIME_DELTA (secs)"
% end_fsp_run

\subsection{Boot compute nodes}

At this point, the {\em master} server should be able to boot the newly defined
compute nodes. Assuming that the compute node BIOS settings are configured to
boot over PXE, all that is required to initiate the provisioning process is to
power cycle each of the desired hosts using IPMI access.
The following commands use the \texttt{ipmitool} utility to initiate power
resets on each of the four compute hosts. Note that the utility requires that
the \texttt{IPMI\_PASSWORD} environment variable be set with the local BMC password in
order to work interactively.

% begin_ohpc_run
% ohpc_comment_header Boot compute nodes \ref{sec:boot_computes}
\begin{lstlisting}[language=bash,keywords={},upquote=true]
[sms](*\#*) for ((i=0; i<${num_computes}; i++)) ; do
              ipmitool -E -I lanplus -H ${c_bmc[$i]} -U ${bmc_username} chassis power reset
        done
\end{lstlisting} 
% end_ohpc_run

Once kicked off, the boot process should take less than 5 minutes (depending on
BIOS post times) and you can verify that the compute hosts are available via
ssh, or via parallel ssh tools to multiple hosts. For example, to run a command
on the newly imaged compute hosts using \texttt{pdsh}, execute the following:

%\iftoggle{isCentOS}{\clearpage}
  
\begin{lstlisting}[language=bash]
[sms](*\#*) pdsh -w c[1-4] uptime
c1  05:03am  up   0:02,  0 users,  load average: 0.20, 0.13, 0.05
c2  05:03am  up   0:02,  0 users,  load average: 0.20, 0.14, 0.06
c3  05:03am  up   0:02,  0 users,  load average: 0.19, 0.15, 0.06
c4  05:03am  up   0:02,  0 users,  load average: 0.15, 0.12, 0.05
\end{lstlisting}
 
The following commands use the \texttt{ipmitool} utility to initiate power
resets on each of the four compute hosts. Note that the utility requires that
the \texttt{IPMI\_PASSWORD} environment variable be set with the local BMC password in
order to work interactively.

% begin_fsp_run
% fsp_validation_comment Boot compute nodes

\begin{lstlisting}[language=bash,keywords={},upquote=true]
[master]$ ipmitool -E -I lanplus -H <c1_bmc> -U root chassis power reset   # power cycle c1
[master]$ ipmitool -E -I lanplus -H <c2_bmc> -U root chassis power reset   # power cycle c2
[master]$ ipmitool -E -I lanplus -H <c3_bmc> -U root chassis power reset   # power cycle c3
[master]$ ipmitool -E -I lanplus -H <c4_bmc> -U root chassis power reset   # power cycle c4
\end{lstlisting} 

% begin_fsp_run 
% fsp_validation_comment - optionally add 4 more hosts
% [master]$ if [[ $USE_DYNAMIC_HOSTNUM -eq 1 && $NUM_COMPUTES -ge 8 ]];then
% [master]$    ipmitool -E -I lanplus -H <c5_bmc> -U root chassis power reset
% [master]$    ipmitool -E -I lanplus -H <c6_bmc> -U root chassis power reset
% [master]$    ipmitool -E -I lanplus -H <c7_bmc> -U root chassis power reset
% [master]$    ipmitool -E -I lanplus -H <c8_bmc> -U root chassis power reset
% [master]$ fi
% end_fsp_run

% begin_fsp_run 
% fsp_validation_comment - optionally add 8 more hosts
% [master]$ if [[ $USE_DYNAMIC_HOSTNUM -eq 1 && $NUM_COMPUTES -ge 8 ]];then
% [master]$    ipmitool -E -I lanplus -H <c9_bmc> -U root chassis power reset
% [master]$    ipmitool -E -I lanplus -H <c10_bmc> -U root chassis power reset
% [master]$    ipmitool -E -I lanplus -H <c11_bmc> -U root chassis power reset
% [master]$    ipmitool -E -I lanplus -H <c12_bmc> -U root chassis power reset
% [master]$    ipmitool -E -I lanplus -H <c13_bmc> -U root chassis power reset
% [master]$    ipmitool -E -I lanplus -H <c14_bmc> -U root chassis power reset
% [master]$    ipmitool -E -I lanplus -H <c15_bmc> -U root chassis power reset
% [master]$    ipmitool -E -I lanplus -H <c16_bmc> -U root chassis power reset
% [master]$ fi
% end_fsp_run

% end_fsp_run

Once kicked off, the boot process should take less than 2 minutes and you can
verify that the compute hosts are available via ssh, or via parallel ssh tools to multiple
hosts. For example, to wait 120 seconds, and run a command on the newly imaged
compute hosts using \texttt{pdsh}, execute the following:

% begin_fsp_run
\begin{lstlisting}[language=bash]
[master]$ sleep 120
[master]$ /opt/fsp/admin/pdsh/bin/pdsh -w c[1-4] uptime
c1  05:03am  up   0:02,  0 users,  load average: 0.20, 0.13, 0.05
c2  05:03am  up   0:02,  0 users,  load average: 0.20, 0.14, 0.06
c3  05:03am  up   0:02,  0 users,  load average: 0.19, 0.15, 0.06
c4  05:03am  up   0:02,  0 users,  load average: 0.15, 0.12, 0.05
\end{lstlisting}
% end_fsp_run



\section{Install \FSP{} Development Components}

The install procedure outlined in \S\ref{sec:basic_install}
highlighted the steps necessary to install a {\em master} host,
assemble and customize a {\em compute} image, and provision several
compute hosts from bare-metal.  With these steps completed, 
%a {\em compute} Once the completion of the
%basic install procedure outlined in Section~\ref{sec:basic_install} is
%complete, 
additional \FSP{}-provided packages can now be added to support a flexible HPC
development environment including development tools, C/C++/Fortran compilers,
MPI stacks, and a variety of 3rd party libraries. The following subsections
highlight the additional software installation procedures, including the
addition of Intel licensed software (e.g. Composer compiler suite, \Intel{}
MPI). It is assumed that the end-site administrator will procure and install
the necessary licenses in order to use the Intel proprietary software.

\subsection{Development Tools}
To aid in general development efforts, \OHPC{} provides recent versions of the \GNU{}
autotools collection, the Valgrind memory debugger, EasyBuild, and R. These can be installed as follows:

% begin_ohpc_run
% ohpc_comment_header Install Development Tools \ref{sec:install_dev_tools}
\begin{lstlisting}[language=bash,keywords={},literate={-}{-}1]
[sms](*\#*) (*\groupinstall*) ohpc-autotools
[sms](*\#*) (*\install*) valgrind-ohpc
[sms](*\#*) (*\install*) EasyBuild-ohpc
[sms](*\#*) (*\install*) spack-ohpc
[sms](*\#*) (*\install*) R_base-ohpc            
\end{lstlisting}
% end_ohpc_run


\subsection{Compilers}
\FSP{} presently packages two compiler families ({\GNU{}} and {\Intel{}
  Parallel Studio}) that are integrated within the underlying
modules-environment system in a hierarchical fashion. End users of a \FSP{}
system can choose to access one compiler at a time and will be presented with
additional compiler-dependent software as a function of which compiler
toolchain is currently loaded. Each compiler toolchain can be installed
separately and the following commands illustrate the installation of both along
with any necessary dependencies:

% begin_fsp_run
% fsp_validation_comment Install compilers
\begin{lstlisting}[language=bash]
[master]$ (*\install*) gnu-compilers-fsp intel-compilers-devel-fsp
\end{lstlisting}
% end_fsp_run


\subsection{Performance Tools}
\OHPC{} provides a variety of open-source tools to aid in application 
performance analysis (refer to Appendix~\ref{appendix:manifest} for a listing
of available packages). This group of tools can be installed as follows:

% begin_ohpc_run
% ohpc_comment_header Install Performance Tools \ref{sec:install_perf_tools}
\begin{lstlisting}[language=bash,keywords={},literate={-}{-}1]
[sms](*\#*) (*\groupinstall*) ohpc-perf-tools-gnu
\end{lstlisting}
% end_ohpc_run




\subsection{MPI Stacks} \label{sec:mpi}
For MPI development support, \OHPC{} presently provides pre-packaged builds for
three MPI families: 

\begin{itemize*}
\item \Intel{}~MPI
\item OpenMPI
\item MVAPICH2
\end{itemize*}
 For ABI consistency, each of the open-source MPI families (OpenMPI and
 MVAPICH2) is built against each of the two supported compiler families
 resulting in total of four build combinations.  The \Intel{} MPI stack is also
 configured to support both the \GNU{} and Intel compiler toolchain directly, but
 is packaged as a single RPM. Installation of all of the MPI family instances,
 can be accomplished via the following command. Note the use of wildcards
 (\texttt{*}) in this example in order to install both \GNU{} and Intel builds for
 OpenMPI and MVAPICH2. In addition to installing available MPI toolchains,
 the example below includes installation of a suite of pre-built MPI
 benchmarks (IMB) that are also provided via \OHPC{}.

% begin_ohpc_run
% fsp_comment_header Install MPI Stacks \ref{sec:mpi}
\begin{lstlisting}[language=bash]
[master]$ (*\install*)  openmpi-*-ohpc mvapich2-*-ohpc intel-mpi-ohpc
[master]$ (*\groupinstall*) ohpc-imb
\end{lstlisting}
% end_ohpc_run


% begin_fsp_run
% fsp_validation_comment Install licenses
% \begin{lstlisting}[language=bash,keywords={},upquote=true]
% [master]$ (*\install*) intel_licenses
% \end{lstlisting}
% end_fsp_run

\subsection{3rd Party Libraries and Tools} \label{sec:3rdparty}

\FSP{} provides pre-packaged builds for a number of popular open-source
tools and libraries used by HPC applications and developers. For
example, \FSP{} provides builds for \FFTW{} and \hdffive{} (including serial and parallel
I/O support), and the \GNU{} Scientific Library (GSL). Again, multiple builds of
each package are available in the \FSP{} repository to support multiple compiler
and MPI family combinations where appropriate. The general naming convention
for builds provided by \FSP{} is to append the compiler and MPI family name that
the library was built against directly into the package name. For example,
libraries that do not require MPI as part of the build process adopt the
following RPM name: \\

\noindent
\texttt{package-<compiler\_family>-fsp-<package\_version>-<release>.rpm} \\

\noindent Packages that require MPI as part of the build expand upon this convention to
additionally include the MPI family name as follows: \\

\noindent
\texttt{package-<compiler\_family>-<mpi\_family>-fsp-<package\_version>-<release>.rpm} \\

To illustrate this further, the command below queries the locally configured
repositories to identify all of the available \FFTW{} packages that were built
with the \GNU{} toolchain. The resulting output that is included shows that
pre-built versions are available for each of the supported MPI families
presented in \S\ref{sec:mpi}.

\begin{lstlisting}[language=bash]
[master]$ yum search fftw-gnu fsp
Loaded plugins: fastestmirror
Loading mirror speeds from cached hostfile
=========================== N/S matched: fftw-gnu, fsp ===========================
fftw-gnu-impi-fsp.x86_64 : A Fast Fourier Transform library
fftw-gnu-mvapich2-fsp.x86_64 : A Fast Fourier Transform library
fftw-gnu-openmpi-fsp.x86_64 : A Fast Fourier Transform library
\end{lstlisting}

Note that \FSP{} provided 3rd party builds are configured to be installed
into a common top-level repository so that they can be easily exported to
desired hosts within the cluster. This common top-level path
(\texttt{/opt/fsp/pub}) was previously configured to be mounted on {\em
  compute} nodes in \S\ref{sec:master_customization}, so the packages will be
immediately available for use on the cluster after installation on the {\em
  master} host.  For convenience, \FSP{} provides package aliases for these 3rd
party libraries and utilities that can be used to install all of the available
compiler/MPI family permutations. To install all of the available package
offerings within \FSP{}, issue the following:

\begin{center}
\begin{tcolorbox}[]
\small
\OHPC{}-provided 3rd party builds are configured to be installed
into a common top-level repository so that they can be easily exported to
desired hosts within the cluster. This common top-level path
(\path{/opt/ohpc/pub}) was previously configured to be mounted on {\em
 compute} nodes in \S\ref{sec:master_customization}, so the packages will be
immediately available for use on the cluster after installation on the {\em
 master} host.
\end{tcolorbox}
\end{center}

%\iftoggle{isCentOS}{\clearpage}
%\nottoggle{isCentOS}{\clearpage}

For convenience, \OHPC{} provides package aliases for these 3rd party libraries
and utilities that can be used to install available libraries for use with the
GNU compiler family toolchain. For parallel libraries, aliases are grouped by
MPI family toolchain so that administrators can choose a subset should they
favor a particular MPI stack.  Please refer to Appendix~\ref{appendix:manifest}
for a more detailed listing of all available packages in each of these functional
areas. To install all available package offerings within \OHPC{}, issue the
following:

% begin_ohpc_run
% ohpc_comment_header Install 3rd Party Libraries and Tools \ref{sec:3rdparty}
\begin{lstlisting}[language=bash,keywords={},upquote=true,keepspaces]
# Install 3rd party libraries/tools meta-packages built with GNU toolchain
[sms](*\#*) (*\install*) ohpc-gnu9-serial-libs
[sms](*\#*) (*\install*) ohpc-gnu9-io-libs
[sms](*\#*) (*\install*) ohpc-gnu9-python-libs
[sms](*\#*) (*\install*) ohpc-gnu9-runtimes
\end{lstlisting}
% end_ohpc_run






% begin_fsp_run 
% fsp_validation_comment - time update
% [master]$ TIME_T2=$(date +%s) ; TIME_DELTA=$(($TIME_T2-$TIME_T1)) ; TIME_T1=$TIME_T2
% [master]$ echo " "; echo "END TIMING = $TIME_DELTA (secs)"
% end_fsp_run

\vspace*{0.2cm}

\section{Resource Manager Startup}
In section \S\ref{sec:basic_install}, the \SLURM{} resource manager was installed
and configured for use on both the {\em master} host and {\em compute} node
instances. With the cluster nodes up and functional, we can now startup the
resource manager services in preparation for running user jobs. Generally, this
is a two-step process that requires starting up the controller daemons on the {\em
  master} host and the client daemons on each of the {\em compute} hosts.  
%Since the {\em compute} hosts were booted into an image that included the SLURM client
%components, the daemons should already be running on the {\em compute}
%hosts. 
Note that \SLURM{} leverages the use of the {\em munge} library to provide
authentication services and this daemon also needs to be running on all hosts
within the resource management pool. 
%The munge daemons should already
%be running on the {\em compute} subsystem at this point, 
The following commands can be used to startup the necessary services to support
resource management under \SLURM{}.

\vspace*{0.3cm}

% begin_fsp_run
% fsp_comment_header Resource Manager Startup \ref{sec:rms_startup}
\begin{lstlisting}[language=bash]
# start munge and slurm controller on master host
[master]$ service munge start
[master]$ service slurmctld start

#  start munge and slurm clients on compute hosts
[master]$ pdsh -w c[1-4] service munge start
[master]$ pdsh -w c[1-4] service slurmd start
\end{lstlisting}
% end_fsp_run

In the default configuration, the {\em compute} hosts will be initialized in an
{\em unknown} state. To place the hosts into production such that they are
eligible to schedule user jobs, issue the following:

% begin_fsp_run
\begin{lstlisting}[language=bash]
[master]$ scontrol update nodename=c[1-4] state=idle
\end{lstlisting}
% end_fsp_run


% Additional recipe commands for additional computes.

% begin_fsp_run
% fsp_validation_newline
% fsp_validation_comment Update basic slurm config if additional computes defined
% fsp_command if [ ${num_computes} -gt 4 ];then
% fsp_command    perl -pi -e "s/^NodeName=(\S+)/NodeName=c[1-${num_computes}]/" /etc/slurm/slurm.conf
% fsp_command    perl -pi -e "s/^PartitionName=normal Nodes=(\S+)/PartitionName=noarmal Nodes=c[1-${num_computes}]/" /etc/slurm/slurm.conf
% fsp_command    service slurmctld restart
% fsp_command    pdsh -w c[5-$num_computes] service munge start
% fsp_command    pdsh -w c[5-$num_computes] service slurmd start
% fsp_command fi
% end_fsp_run


\section{Run a Test Job}
With the resource manager enabled for production usage, users should now be
able to run jobs. To demonstrate this, we will add a ``test'' user on the {\em master}
host that can be used to run an example job.

% begin_ohpc_run
\begin{lstlisting}[language=bash,keywords={}]
[sms](*\#*) useradd -m test
\end{lstlisting}
% end_ohpc_run

\Warewulf{} installs a utility on the compute nodes to automatically 
synchronize known files from the provisioning server at five minute intervals. In this
recipe, recall that we previously registered credential files with Warewulf (e.g. passwd,
group, and shadow) so that these files would be propagated during compute node
imaging. However, with the addition of a new ``test'' user above, the files
have been outdated and we need to update the Warewulf database to incorporate
the additions. This re-sync process can be accomplished as follows:

% begin_ohpc_run
\begin{lstlisting}[language=bash,keywords={}]
[sms](*\#*) wwsh file resync passwd shadow group
\end{lstlisting}
% end_ohpc_run

% begin_ohpc_run
% ohpc_command sleep 2
% end_ohpc_run

\begin{center}
\begin{tcolorbox}[]
\small
After re-syncing to notify Warewulf of file modifications made on the {\em
master} host, it should take approximately 5 minutes for the changes to
propagate. However, you can also manually pull the changes from compute nodes
via the following:
% begin_ohpc_run
\begin{lstlisting}[language=bash,keywords={}]
[sms](*\#*) pdsh -w c[1-4] /warewulf/bin/wwgetfiles 
\end{lstlisting}
% end_ohpc_run
\end{tcolorbox}
\end{center}

%We must first update the database with the change we just made to passwd and 
%shadow files, and since we are impatient we will trigger the compute node pull.

\OHPC{} includes a simple ``hello-world'' MPI application in the
\path{/opt/ohpc/pub/examples} directory that can be used for this quick
compilation and execution. \OHPC{} also provides a companion job-launch script
named \texttt{prun} that is installed in concert with the pre-packaged MPI
toolchains. At present, \OHPC{} is unable to include the PMI process management
server normally included within \SLURM{} which implies that \texttt{srun} cannot
be use for MPI job launch. Instead, native job launch mechanisms provided by
the MPI stacks are utilized and \texttt{prun} abstracts this process for the
various stacks to retain a single launch command.

%\iftoggle{isSLES_ww_slurm_x86}{\clearpage}
%\iftoggle{isCentOS_ww_slurm_x86}{\clearpage}

  
\subsection{Interactive execution}
To use the newly created ``test'' account to compile and execute the
application {\em interactively} through the resource manager, execute the
following (note the use of \texttt{prun} for parallel job launch which summarizes
the underlying native job launch mechanism being used):

\begin{lstlisting}[language=bash,keywords={}]
# Switch to "test" user
[sms](*\#*) su - test

# Compile MPI "hello world" example
[test@sms ~]$ mpicc -O3 /opt/ohpc/pub/examples/mpi/hello.c

# Submit interactive job request and use prun to launch executable
[test@sms ~]$ srun -n 8 -N 2 --pty /bin/bash

[test@c1 ~]$ prun ./a.out

[prun] Master compute host = c1
[prun] Resource manager = slurm
[prun] Launch cmd = mpiexec.hydra -bootstrap slurm ./a.out

 Hello, world (8 procs total)
    --> Process #   0 of   8 is alive. -> c1
    --> Process #   4 of   8 is alive. -> c2
    --> Process #   1 of   8 is alive. -> c1
    --> Process #   5 of   8 is alive. -> c2
    --> Process #   2 of   8 is alive. -> c1
    --> Process #   6 of   8 is alive. -> c2
    --> Process #   3 of   8 is alive. -> c1
    --> Process #   7 of   8 is alive. -> c2
\end{lstlisting}

\begin{center}
\begin{tcolorbox}[]
The following table provides approximate command equivalences between SLURM and
PBS Pro:

\vspace*{0.15cm}
\input common/rms_equivalence_table
\end{tcolorbox}
\end{center}
\nottoggle{isCentOS}{\clearpage}

\iftoggle{isCentOS}{\clearpage}

\subsection{Batch execution}

For batch execution, \OHPC{} provides a simple job script for reference (also
housed in the \path{/opt/ohpc/pub/examples} directory. This example script can
be used as a starting point for submitting batch jobs to the resource manager
and the example below illustrates use of the script to submit a batch job for
execution using the same executable referenced in the previous interactive example.

\begin{lstlisting}[language=bash,keywords={}]
# Copy example job script
[test@sms ~]$ cp /opt/ohpc/pub/examples/slurm/job.mpi .

# Examine contents (and edit to set desired job sizing characteristics)
[test@sms ~]$ cat job.mpi
#!/bin/bash

#SBATCH -J test               # Job name
#SBATCH -o job.%j.out         # Name of stdout output file (%j expands to %jobId)
#SBATCH -N 2                  # Total number of nodes requested
#SBATCH -n 16                 # Total number of mpi tasks #requested
#SBATCH -t 01:30:00           # Run time (hh:mm:ss) - 1.5 hours

# Launch MPI-based executable

prun ./a.out

# Submit job for batch execution
[test@sms ~]$ sbatch job.mpi
Submitted batch job 339
\end{lstlisting}

\begin{center}
\begin{tcolorbox}[]
\small
The use of the \texttt{\%j} option in the example batch job script shown is a convenient
way to track application output on an individual job basis. The \texttt{\%j} token
is replaced with the \SLURM{} job allocation number once assigned (job~\#339 in
this example).
\end{tcolorbox}
\end{center}




%\newcommand{\firstColWidth}{2.3cm}
%\newcommand{\secondColWidth}{1.25cm}

\clearpage

\definecolor{Gray}{gray}{0.5}
\newcommand{\captionSpace}{-0.15cm}
\newcommand{\tabSpaceBot}{1.0cm}
\captionsetup{justification=raggedright,singlelinecheck=false}

\subsection{Package Manifest} \label {appendix:manifest}

\vspace*{0.25cm}
This appendix provides a summary of available meta-package groupings and all of
the individual RPM packages that are available as part of this \OHPC{}
release. The meta-packages provide a mechanism to group related collections of
RPMs by functionality and provide a convenience mechanism for installation.  A
list of the available meta-packages and a brief description is presented in
Table~\ref{table:groups}.

\vspace*{1.25cm}
\begin{table}[h] 
\caption{\bf Available \OHPC{} Meta-packages} \vspace*{\captionSpace{}}
\label{table:groups}
\input manifest/patterns
\end{table}

\newpage
What follows next in this Appendix is a series of tables that summarize the
underlying RPM packages available in this \OHPC{} release. These packages are
organized by groupings based on their general functionality and each table
provides information for the specific RPM name, version, brief summary, and the
web URL where additional information can be obtained for the component. Note
that many of the 3rd party community libraries that are pre-packaged
with \OHPC{} are built using multiple compiler and MPI families. In these cases,
the RPM package name includes delimiters identifying the development
environment for which each package build is targeted.  Additional information
on the \OHPC{} package naming scheme is presented in \S\ref{sec:3rdparty}. 
The relevant package groupings and associated Table references are as follows:

\vspace*{0.1cm}

\begin{itemize*}
\item Administrative tools (Table~\ref{table:admin})
\item Provisioning (Table~\ref{table:provisioning})
\item Resource management (Table~\ref{table:rms})
\item Compiler families (Table~\ref{table:compiler-families})
\item MPI families (Table~\ref{table:mpi-families})
\item Development tools (Table~\ref{table:dev-tools})
\item Performance analysis tools (Table~\ref{table:perf-tools})
\item Distro support packages and dependencies (Table~\ref{table:distro-packages})
\item IO Libraries (Table~\ref{table:io-libs})
\item Runtimes (Table~\ref{table:runtimes})
\item Serial Libraries (Table~\ref{table:serial-libs})
\item Parallel Libraries (Table~\ref{table:parallel-libs})
\end{itemize*}


\newcommand{\firstColWidth}{5.0cm}
\newcommand{\secondColWidth}{1.25cm}

\vspace*{1.0cm}
\urlstyle{same}

% Administration Tools 
\begin{table}[h]
\caption{\bf Administrative Tools} \vspace*{\captionSpace{}} \label{table:admin}
\input manifest/admin
\end{table}
\vspace*{0.5cm}

\renewcommand{\firstColWidth}{6.25cm}
\renewcommand{\secondColWidth}{0.95cm}

% Provisioning
\iftoggleverb{isWarewulf}
\begin{table}[h!]
\caption{\bf Provisioning} \vspace*{\captionSpace{}} \label{table:provisioning}
\input manifest/provisioning
\vspace*{\tabSpaceBot{}}
\end{table}
\fi

\renewcommand{\firstColWidth}{4.1cm}
\renewcommand{\secondColWidth}{1.8cm}

% Resource Management
\begin{table}[h!]
\caption{\bf Resource Management} \vspace*{\captionSpace{}} \label{table:rms}
\input manifest/rms
\vspace*{\tabSpaceBot{}}
\end{table}

\renewcommand{\firstColWidth}{4.7cm}

% Compiler Families
\begin{table}[h!]
\caption{\bf Compiler Families} \vspace*{\captionSpace{}} \label{table:compiler-families}
\input manifest/compiler-families
\vspace*{\tabSpaceBot{}}
\end{table}

\renewcommand{\firstColWidth}{4.7cm}

% MPI Families
\begin{table}[h!]
\caption{\bf MPI Families} \vspace*{\captionSpace{}} \label{table:mpi-families}
\input manifest/mpi-families
\vspace*{\tabSpaceBot{}}
\end{table}

\renewcommand{\firstColWidth}{5.6cm}
\renewcommand{\secondColWidth}{1.5cm}

% Development Tools
\begin{table}[h!]
\caption{\bf Development Tools} \vspace*{\captionSpace{}} \label{table:dev-tools}
\input manifest/dev-tools
\vspace*{\tabSpaceBot{}}
\end{table}

% Development Tools (2)
\begin{table}[h!]
\caption*{Table~\ref{table:dev-tools} (cont): {\bf Development Tools} \vspace*{\captionSpace{}} }
\input manifest/dev-tools2
\vspace*{\tabSpaceBot{}}
\end{table}

\renewcommand{\firstColWidth}{4.5cm}
\renewcommand{\secondColWidth}{1.7cm}

% Perf Tools (1)
\begin{table}[h!]
\caption{\bf Performance Analysis Tools} \vspace*{\captionSpace{}} \label{table:perf-tools}
\input manifest/perf-tools
\vspace*{\tabSpaceBot{}}
\end{table}

% Perf Tools (2)
\begin{table}[h!]
\caption*{Table~\ref{table:perf-tools} (cont): {\bf Performance Analysis Tools} \vspace*{\captionSpace{}} }
\input manifest/perf-tools2
\vspace*{\tabSpaceBot{}}
\end{table}

% Perf Tools (3)
\begin{table}[h!]
\caption*{Table~\ref{table:perf-tools} (cont): {\bf Performance Analysis Tools} \vspace*{\captionSpace{}} }
\input manifest/perf-tools3
\vspace*{\tabSpaceBot{}}
\end{table}

% Distro Packages
\begin{table}[h!]
\caption{\bf Distro Support Packages/Dependencies} \vspace*{\captionSpace{}} \label{table:distro-packages}
\input manifest/distro-packages
\vspace*{\tabSpaceBot{}}
\end{table}

\renewcommand{\firstColWidth}{5.1cm}
\renewcommand{\secondColWidth}{1.4cm}

% Lustre
\begin{table}[h!]
\caption{\bf Lustre} \vspace*{\captionSpace{}} \label{table:lustre}
\input manifest/lustre
\vspace*{\tabSpaceBot{}}
\end{table}

% IO Libs (1)
\begin{table}[h!]
\caption{\bf IO Libraries} \vspace*{\captionSpace{}} \label{table:io-libs}
\input manifest/io-libs
\vspace*{\tabSpaceBot{}}
\end{table}

% IO Libs (2)
\begin{table}[h!]
\caption*{Table~\ref{table:io-libs} (cont): {\bf IO Libraries} \vspace*{\captionSpace{}} }
\input manifest/io-libs2
\vspace*{\tabSpaceBot{}}
\end{table}

% IO Libs (3)
\begin{table}[h!]
\caption*{Table~\ref{table:io-libs} (cont): {\bf IO Libraries} \vspace*{\captionSpace{}} }
\input manifest/io-libs3
\vspace*{\tabSpaceBot{}}
\end{table}

%\renewcommand{\firstColWidth}{4.5cm}
%\renewcommand{\secondColWidth}{1.5cm}

% Runtimes
\clearpage
\begin{table}[h!]
\caption{\bf Runtimes} \vspace*{\captionSpace{}} \label{table:runtimes}
\input manifest/runtimes
\vspace*{\tabSpaceBot{}}
\end{table}

% Serial libs
\begin{table}[h!]
\caption{\bf Serial/Threaded Libraries} \vspace*{\captionSpace{}} \label{table:serial-libs}
\input manifest/serial-libs
\vspace*{\tabSpaceBot{}}
\end{table}

% Parallel libs (1)
\begin{table}[h!]
\caption{\bf Parallel Libraries} \vspace*{\captionSpace{}} \label{table:parallel-libs}
\input manifest/parallel-libs
\vspace*{\tabSpaceBot{}}
\end{table}

% Parallel libs (2)
\begin{table}[h!]
\caption*{Table~\ref{table:parallel-libs} (cont): {\bf Parallel Libraries} \vspace*{\captionSpace{}} }
\input manifest/parallel-libs2
\vspace*{\tabSpaceBot{}}
\end{table}

% Parallel libs (3)
\begin{table}[h!]
\caption*{Table~\ref{table:parallel-libs} (cont): {\bf Parallel Libraries} \vspace*{\captionSpace{}} }
\input manifest/parallel-libs3
\vspace*{\tabSpaceBot{}}
\end{table}

% Parallel libs (4)
\begin{table}[h!]
\caption*{Table~\ref{table:parallel-libs} (cont): {\bf Parallel Libraries} \vspace*{\captionSpace{}} }
\input manifest/parallel-libs4
\vspace*{\tabSpaceBot{}}
\end{table}







% begin_fsp_run 
% fsp_validation_comment - time update
% [master]$ TIME_T2=$(date +%s) ; TIME_DELTA=$(($TIME_T2-$TIME_T1)) ; TIME_T1=$TIME_T2
% [master]$ echo " "; echo "END TIMING = $TIME_DELTA (secs)"
% end_fsp_run

\end{document}

